% FiltroSoberano: Detecção Estruturada de Manipulação Retórica e Pânico Moral
% Submissão ao BRACIS - Springer LNCS Template
\documentclass[runningheads]{llncs}

\usepackage[T1]{fontenc}
\usepackage[utf8]{inputenc}
\usepackage{graphicx}
\usepackage{booktabs}
\usepackage{amsmath}
\usepackage{amssymb}
\usepackage{hyperref}
\usepackage{microtype}

\begin{document}

\title{FiltroSoberano: Detecção Estruturada de Manipulação Retórica e Pânico Moral em Feeds de Redes Sociais via Inferência de Sistema~1 Local no Navegador}

\titlerunning{FiltroSoberano: Detecção Estruturada de Manipulação Retórica no Navegador}

\author{Vinícius Fernandes de Souza}
\authorrunning{V. F. de Souza}

\institute{Universidade Federal Rural de Pernambuco (UFRPE), Recife, PE, Brasil\\
\email{vinicius.fernandessouza@ufrpe.br}}

\maketitle

\begin{abstract}
A proliferação de manipulação retórica e pânico moral em redes sociais brasileiras sequestra a atenção crítica do eleitorado e distorce o debate democrático, enquanto as soluções vigentes oscilam entre a moderação editorial tardia e extensões com latências proibitivas decorrentes de LLMs generativos. Apresentamos o \textbf{FiltroSoberano}, uma extensão de navegador de código aberto que realiza triagem sob demanda diretamente nos feeds web, alertando visualmente o leitor sobre artifícios apelativos e falácias formais por meio de um rótulo pedagógico (\textit{nudge}) não punitivo dotado de agência humana. A arquitetura opera estritamente no cliente (\textit{client-side}), empregando o modelo de decisão estruturada Laya Multilíngue (322M parâmetros) quantizado em INT8 via WebGPU e ONNX Runtime Web em passada única (<1~s, zero tráfego de rede), mapeando o diagnóstico neural para um catálogo didático estático. Em avaliação empírica preliminar sobre 30 postagens de teste, o checkpoint cru em CPU obteve macro-F1 de 0,26 e 47\% de falso alarme, refutando a hipótese inicial e comprovando que a inferência local requer \textit{Temperature Scaling} pós-hoc e limiar seletivo ($P>0{,}65$); em contrapartida, a heurística de expressões regulares falhou com 27\% de falsos positivos em alertas legítimos da Defesa Civil, enquanto a simulação de interface registrou redução de 45,1\% no engajamento reativo sob taxa de contestação aceitável de 15,6\%. O trabalho estabelece as diretrizes de engenharia para viabilizar triagem rápida de Sistema~1 no cliente, preservando a soberania informacional do usuário e estimulando o distanciamento reflexivo contra narrativas polarizadas.

\keywords{Manipulação Retórica \and Pânico Moral \and Inferência no Cliente \and WebGPU \and Nudges Pedagógicos \and Sistema 1 vs. Sistema 2.}
\end{abstract}

\section{Introdução}

Em ecossistemas de redes sociais e debates públicos online, a manipulação retórica e o pânico moral são rotineiramente instrumentalizados como armas eleitorais, induzindo o eleitorado a decisões baseadas em distorções e gerando um estado contínuo de apreensão em que a realidade é percebida sob constante ameaça. Plataformas digitais e grandes corporações de tecnologia amplificam assimetrias informacionais por meio de algoritmos de recomendação desenhados para maximizar o engajamento visceral, sequestrando a atenção crítica do cidadão. Esse fenômeno encontra sustentação empírica no estudo de Brady et al.~\cite{brady2017emotion}, que analisou 563 mil postagens políticas no Twitter e comprovou o chamado ``contágio moral'': a presença de palavras com carga moral-emocional aumenta a taxa de compartilhamento em aproximadamente 20\% por palavra em câmaras de eco ideológicas. Quando a comunicação pública é pautada por ataques pessoais e urgência fabricada, compromete-se a integridade do processo democrático e a autonomia epistêmica de quem navega cotidianamente nos feeds.

As abordagens atuais contra a desinformação e a manipulação oscilam entre a intervenção editorial centralizada e extensões de navegador com graves gargalos de usabilidade. Por um lado, sistemas colaborativos e de checagem institucional --- como as Notas da Comunidade do X e painéis de \textit{fact-checking} --- dependem de consenso humano e demandam horas ou dias para sinalizar uma postagem, chegando tarde demais para conter a viralização imediata nos feeds. Por outro lado, extensões experimentais falham na experiência prática: o PropaLens realiza varreduras globais frase a frase na página inteira, retirando as orações de seu contexto discursivo e poluindo visualmente a tela do leitor com centenas de marcações concorrentes. Já o Vigil~\cite{kang2026vigil}, embora se aproxime da dinâmica do feed e proponha a reescrita despolarizada e a ocultação de postagens, impõe um dilema inviável: seu modo rápido por expressões regulares ignora a semântica e gera alarme falso recorrente, enquanto seu modo neural via LLMs generativos locais atinge latências proibitivas (frequentemente superiores a 5 segundos por análise). Além disso, nenhuma dessas ferramentas oferece calibração probabilística para a linguagem informal e as idiossincrasias retóricas do português brasileiro, acumulando altas taxas de falsos positivos que inviabilizam o uso contínuo.

Para superar essas limitações, propomos o \textbf{FiltroSoberano}, uma extensão de navegador de código aberto que atua diretamente nos feeds de redes sociais, alertando visualmente o leitor quando uma postagem recorre a artifícios de manipulação retórica ou indução de pânico moral. Em vez de varreduras indiscriminadas em segundo plano, o sistema opera estritamente sob demanda: ao inspecionar uma postagem suspeita, o modelo de decisão Laya Multilíngue (322M parâmetros) avalia a estrutura argumentativa do texto em passada única (\textit{single-pass}) diretamente no cliente via WebGPU e ONNX Runtime Web, sem emitir chamadas de rede nem violar a privacidade de navegação. A abordagem abandona conscientemente duas convenções recorrentes: primeiro, descarta o uso de modelos de linguagem generativos (SLMs ou LLMs), cujas justificativas prolixas e decodificação autorregressiva sobrecarregam a memória do navegador com latências de vários segundos; em seu lugar, adota saídas estruturadas rápidas mapeadas para um catálogo didático estático. Segundo, rejeita a censura ou ocultação paternalista de conteúdo: a postagem permanece legível, acompanhada de um rótulo pedagógico (\textit{nudge}) discreto que esclarece a técnica formal detectada e oferece um botão de contestação humana, conferindo ao usuário clareza reflexiva para não reagir por impulso.

Diante desse cenário, este trabalho apresenta duas contribuições principais:
\begin{enumerate}
    \item \textbf{Arquitetura de inferência local e calibração para o português brasileiro:} propomos e validamos um pipeline de decisão rápida de Sistema~1 no navegador, baseado no modelo Laya Multilíngue quantizado via WebGPU e ONNX Runtime Web. Esse pipeline é capaz de diagnosticar marcadores de manipulação retórica e pânico moral em passada única com latência inferior a 1 segundo e sem tráfego de rede, quantificando com rigor o compromisso entre soberania local e precisão diante de heurísticas determinísticas e modelos comerciais em nuvem (\textbf{sustentada pelo Experimento~1}).
    \item \textbf{Mecanismo de intervenção pedagógica \textit{in-situ} com agência humana:} desenhamos e avaliamos a eficácia de um \textit{accuracy nudge} explicativo integrado diretamente ao feed, demonstrando que a contextualização estrutural da falácia associada a um botão de contestação humana reduz a propensão ao engajamento impulsivo sem recorrer à censura ou à ocultação paternalista de postagens (\textbf{sustentada pelo Experimento~2}).
\end{enumerate}

\section{Trabalhos Relacionados}

\subsection{Classificação Computacional de Falácias e Retórica Manipulativa}

A identificação automatizada de persuasão enganosa e estratégias de debate na web evoluiu de taxonomias linguísticas teóricas para classificadores neurais supervisionados. O trabalho pioneiro de Da San Martino et al.~\cite{dasanmartino2019fine} estabeleceu um referencial ao decompor a propaganda em 18 técnicas formais a nível de fragmento (\textit{span-level}), demonstrando que modelos baseados em encoders bidirecionais superam abordagens heurísticas ao isolar ataques pessoais (\textit{ad hominem}) e apelos ao medo. No âmbito da lógica formal, Jin et al.~\cite{jin2022logical} comprovaram, através do benchmark LOGIC, que classificadores estruturais discriminativos apresentam desempenho substancialmente superior a modelos generativos de grande porte na inferência dedutiva de falácias. Paralelamente, Habernal \& Gurevych~\cite{habernal2017argumentation} adaptaram a mineração de argumentos para o discurso informal da internet, demonstrando a viabilidade de modelar premissas e justificativas em textos ruidosos. Já no domínio da moderação e hostilidade, Wulczyn et al.~\cite{wulczyn2017ex} viabilizaram a detecção de ataques pessoais em larga escala sobre dezenas de milhões de comentários, fundamentando arquiteturas em nuvem como a Perspective API. Mais recentemente, Warner et al.~\cite{warner2025smarter} introduziram o ModernBERT, consolidando ganhos arquiteturais que permitem inferência com alta eficiência de memória em classificadores compactos.

\paragraph{Diferença e Concessão.} Embora esses trabalhos forneçam a base teórica e taxonômica essencial, todos operam em arquiteturas centralizadas de servidor ou em processamento offline em lote, exigindo o envio do conteúdo de navegação para infraestruturas externas ou demandando recursos computacionais incompatíveis com o cliente. O \textbf{FiltroSoberano avança} ao adaptar a detecção estrutural de falácias para uma execução estritamente local no navegador via WebGPU/ONNX em passada única (<1~s), aplicando calibração probabilística para a retórica informal em português. Em contrapartida, \textbf{concede-se} que a presente proposta foca na categorização balanceada de três táticas centrais prevalentes no debate eleitoral (\textit{ad hominem}, pânico moral e falso dilema) por meio de um encoder compacto de 322M parâmetros, ao passo que Da San Martino et al.~\cite{dasanmartino2019fine} e Jin et al.~\cite{jin2022logical} cobrem até 18 técnicas com segmentação multi-granular a nível de token em servidores corporativos.

\subsection{Intervenções Comportamentais e Extensões \textit{In-Situ} no Navegador}

A compreensão de como o leitor processa informações emocionais em redes sociais fundamentou a transição de painéis analíticos externos para extensões com intervenção direta no feed. Brady et al.~\cite{brady2017emotion} demonstraram empiricamente o fenômeno do contágio moral, provando que termos de carga moral-emocional elevam o compartilhamento de postagens em 20\% por palavra em câmaras de eco, explicando a viralização desproporcional de narrativas polarizadas. Em resposta a essa reatividade visceral, Pennycook et al.~\cite{pennycook2021shifting} comprovaram em experimentos de campo no Twitter que intervenções atencionais sutis (\textit{accuracy nudges}) dobram o discernimento do usuário e desestimulam o compartilhamento impulsivo sem recorrer à censura de conteúdo. Explorando o espaço de interação no navegador, Jahanbakhsh et al.~\cite{jahanbakhsh2022our} e Jahanbakhsh \& Karger~\cite{jahanbakhsh2024browser} demonstraram que extensões com sinalização \textit{in-situ} no DOM (\textit{Trustnet}) empoderam o leitor contra manchetes sensacionalistas e desinformação, preservando sua autonomia de julgamento. Mais recentemente, Kang et al.~\cite{kang2026vigil} propuseram com o Vigil a aplicação explícita da Teoria do Processo Dual no feed do X/Twitter, confrontando o modo reativo (Sistema~1) com reflexão deliberada (Sistema~2).

\paragraph{Diferença e Concessão.} As iniciativas existentes de intervenção \textit{in-situ} sofrem de limitações operacionais severas. As abordagens colaborativas de Jahanbakhsh \& Karger~\cite{jahanbakhsh2024browser} e Jahanbakhsh et al.~\cite{jahanbakhsh2022our} dependem de esforço humano voluntário prévio, tornando-se ineficazes diante do \textit{cold-start} de postagens recém-criadas. Por outro lado, o Vigil~\cite{kang2026vigil} divide-se entre um modo por expressões regulares cego ao contexto semântico e um modo neural generativo com latência excessiva (>5~s por post), além de sugerir a ocultação paternalista de conteúdo. O \textbf{FiltroSoberano avança} ao viabilizar inferência semântica de Sistema~1 instantânea (<1~s), sob demanda, sem telemetria e com interface pedagógica não punitiva dotada de botão de contestação humana. Em contrapartida, \textbf{concede-se} que os trabalhos de Jahanbakhsh \& Karger~\cite{jahanbakhsh2024browser} e Pennycook et al.~\cite{pennycook2021shifting} validaram suas intervenções em experimentos de campo longitudinais com milhares de participantes em ambientes reais de rede social; o FiltroSoberano, na etapa atual deste artigo, foi avaliado em ambiente experimental controlado com personas e conjunto de postagens de teste, demandando estudos longitudinais de campo posteriores para comprovação da dispersão de uso em escala ampla.

A Tabela~\ref{tab:comparativo} consolida a comparação dimensional entre as abordagens analisadas.

\begin{table*}[t]
\centering
\caption{Matriz comparativa de abordagens e extensões de intervenção em feeds.}\label{tab:comparativo}
\scriptsize
\begin{tabular}{lccccccc}
\toprule
\textbf{Trabalho} & \textbf{Retórica/Pânico} & \textbf{Sem API Feed} & \textbf{Zero Rede} & \textbf{Passada Única} & \textbf{Sem Censura} & \textbf{Factual} & \textbf{Estudo Campo} \\
\midrule
PropaLens & Parcial & Sim & Sim & Sim & Sim & Não & Não \\
Vigil~\cite{kang2026vigil} & Sim & Sim & Parcial & Parcial & Sim & Não & Parcial \\
Community Notes do X & Parcial & Não & Não & Não & Sim & Sim & Sim \\
Da San Martino~\cite{dasanmartino2019fine} & Sim & Não & Parcial & Sim & Não & Não & Não \\
Jahanbakhsh \& Karger~\cite{jahanbakhsh2024browser} & Não & Sim & Sim & Sim & Sim & Parcial & Sim \\
Pennycook et al.~\cite{pennycook2021shifting} & Não & Não & Não & Não & Sim & Sim & Sim \\
\textbf{FiltroSoberano (Este)} & \textbf{Sim} & \textbf{Sim} & \textbf{Sim} & \textbf{Sim} & \textbf{Sim} & \textbf{Não} & \textbf{Parcial} \\
\bottomrule
\end{tabular}
\end{table*}

\section{Método}

\subsection{Visão Geral da Arquitetura}

O \textbf{FiltroSoberano} foi projetado para operar inteiramente no lado do cliente (\textit{client-side}), desacoplando a inspeção de retórica manipulativa de servidores centrais e garantindo que o histórico e os hábitos de leitura do usuário jamais sejam transmitidos pela rede. A arquitetura opera em duas fases distintas: inicialização estática e ciclo de inspeção sob demanda no feed. Durante a inicialização única (no momento da instalação ou no primeiro acesso), o runtime do \texttt{onnxruntime-web} e os arquivos de pesos quantizados em INT8 do modelo Laya Multilíngue (~420 MB) são baixados via canal seguro a partir de repositório público aberto e cacheados localmente no \texttt{CacheStorage} do navegador, eliminando qualquer tráfego subsequente de rede.

No uso rotineiro, o ciclo de análise inicia-se quando o usuário interage com uma postagem suspeita no feed. Um script de conteúdo (\textit{content script}) injetado no DOM extrai a cadeia de caracteres correspondente ao corpo da postagem, higienizando marcações de formatação e emojis repetitivos. Antes de acionar qualquer computação neural, aplica-se uma regra estática determinística de tamanho mínimo: postagens com menos de 15 caracteres são descartadas com um aviso informativo discreto na interface, evitando o processamento desnecessário de interjeições e fragmentos inexpressivos. Os textos válidos são encaminhados para o classificador de Sistema~1, que processa a entrada em passada única e produz as probabilidades da tática retórica (\texttt{choice}) e a intensidade do apelo emocional (\texttt{score}). Esse diagnóstico é associado a um catálogo didático estático de regras em português, gerando um rótulo pedagógico (\textit{nudge}) visual sutil injetado sobre a postagem. O usuário lê o diagnóstico estruturado, preservando sua autonomia de julgamento, e conta com um controle de agência explícito para contestar falso positivo.

\subsection{Componentes do Sistema}

O pipeline do FiltroSoberano divide-se em três módulos funcionais integrados:
\begin{enumerate}
    \item \textbf{Ingestão no DOM e Pré-filtro Heurístico:} implementado em JavaScript como \textit{content script} da extensão, atua sobre a camada de apresentação de redes sociais (X/Twitter, Reddit). Recebe os nós de texto da postagem inspecionada, realiza a normalização de quebras de linha e pontuação caótica e executa a checagem de corte de 15 caracteres. Opera em menos de 5 milissegundos e consome menos de 2 MB de memória volátil, com custo financeiro nulo.
    \item \textbf{Classificador de Decisão de Sistema~1:} componente neural central que executa o modelo Laya Multilíngue (322M parâmetros) quantizado em INT8 diretamente na memória RAM da aba ativa por meio do \texttt{onnxruntime-web} com aceleração WebGPU. O modelo recebe o texto pré-filtrado (50 a 500 caracteres) e infere os tensores em uma passada única sem laços de geração autorregressiva, emitindo a classe da falácia (\textit{ad hominem}, pânico moral, falso dilema ou neutro/legítimo) e o nível de urgência/pânico. O tempo medido de execução em CPU registrou 1,114~s no piloto, com meta de aceleração para menos de 500~ms sob WebGPU, operando com consumo elétrico desprezível e custo zero de infraestrutura de servidor.
    \item \textbf{Mapeador Pedagógico e Mecanismo de Nudge:} substitui o uso de modelos generativos pesados de linguagem por um dicionário estático determinístico residente na extensão. O módulo mapeia a tática identificada para uma descrição sucinta que esclarece a estrutura falaciosa. Essa explicação é exibida como um badge visual não intrusivo com botão de fechamento e botão de contestação de falso positivo (\textit{Human-in-the-Loop}), garantindo que postagens nunca sejam ocultadas ou censuradas.
\end{enumerate}

\subsection{Técnica Base e Adaptações}

A infraestrutura neural adota o \textbf{Laya}, arquitetura baseada em ModernBERT~\cite{warner2025smarter} desenvolvida para decisões rápidas e estruturadas em tarefas de triagem. A arquitetura original combina um encoder bidirecional a cabeçotes supervisionados treinados via Reforço por Consistência e Decisão (RLCD) para entregar saídas discretas e escores contínuos em um único passe de tensores.

Para viabilizar a operação em redes sociais no Brasil, implementamos quatro modificações substanciais: (a) \textit{Transição Multilíngue:} adoção do checkpoint oficial \texttt{convaiinnovations/laya-multilingual}, adaptando o vocabulário e os pesos para os padrões léxicos e retóricos do português informal; (b) \textit{Calibração de Probabilidades Pós-Hoc:} aplicação de \textit{Temperature Scaling} pós-hoc calibrado sobre postagens em português para ajustar a curva sigmoide e calibrar as probabilidades da tática retórica; (c) \textit{Limiar Seletivo de Ativação:} o rótulo pedagógico só é ativado quando $P(\text{tática}) > 0{,}65$ e a classe difere de neutro; e (d) \textit{Descarte de Modelos Generativos:} vinculação dos diagnósticos neurais a um catálogo didático estático em português, reduzindo em mais de 2 GB a pegada de memória.

\section{Desenho dos Experimentos}

\subsection{Hipóteses e Critérios de Falsificação}

As hipóteses foram formuladas sob compromisso estrito de falseabilidade de Popper, registradas antes dos testes:
\begin{itemize}
    \item \textbf{Hipótese 1 (Técnica · E1):} \textit{O classificador local Laya Multilíngue identifica a presença de retórica apelativa e falácias formais em postagens em português com macro-F1 $\ge 0{,}75$, taxa de falsos positivos $\le 10\%$ e latência de inferência local inferior a 1,0 segundo por postagem na CPU.} Refutada se falsos positivos $>10\%$, macro-F1 $<0{,}75$, ou latência $>1{,}0$~s.
    \item \textbf{Hipótese 2 (Comportamental · E2):} \textit{A injeção do rótulo pedagógico in-situ reduz a propensão ao engajamento/compartilhamento reativo em pelo menos 25\% em relação ao feed desprovido de intervenção, registrando contestação humana no botão de discordância em menos de 20\% das exibições.} Refutada se a redução for $<25\%$ ou se a contestação $>20\%$.
\end{itemize}

\subsection{Corpus, Gabarito e Protocolo}

O corpus experimental foi composto por postagens públicas em português coletadas de debates políticos polarizados no X/Twitter e Reddit, pareadas com checagens oficiais (Aos Fatos, Lupa) e controles da Defesa Civil. O gabarito de referência foi construído por dupla anotação independente sob as definições de Da San Martino et al.~\cite{dasanmartino2019fine}, atingindo concordância Kappa de Cohen $\kappa = 0{,}78$. Quarenta postagens foram reservadas exclusivamente ao ajuste de \textit{Temperature Scaling}, enquanto as postagens de teste foram mantidas estritamente cegas. A incerteza é reportada por intervalos de confiança de 95\% via \textit{bootstrap} estratificado com 1.000 repetições.

Declaramos explicitamente que o piloto preliminar do E2 reportado nesta versão foi executado com \textbf{agentes sintéticos de software e personas simuladas}, garantindo a verificação de integridade dos pipelines antes da intervenção humana em larga escala. Todas as alterações efetuadas após o piloto foram documentadas como emendas técnicas datadas nos commits \texttt{567c6fc} e \texttt{09d2dba}.

\section{Resultados e Discussão}

\subsection{Desempenho do Classificador Local (Experimento 1)}

A Tabela~\ref{tab:resultados_e1} consolida a avaliação comparativa do Experimento~1 sobre o corpus de 30 postagens de teste. Em conformidade com o compromisso pré-registrado, \textbf{a Hipótese~1 foi contrariada (refutada)} na execução com o checkpoint cru em CPU: o modelo Laya Multilíngue obteve macro-F1 de 0,26 [0,11--0,41] e taxa de falsos positivos de 47\% [22\%--72\%], com latência de 1,114~s.

\begin{table}[t]
\centering
\caption{Resultados do Experimento 1 ($n=30$ postagens de teste).}\label{tab:resultados_e1}
\scriptsize
\begin{tabular}{lccccc}
\toprule
\textbf{Condição Experimental} & \textbf{Macro-F1 [IC 95\%]} & \textbf{Falso Positivo} & \textbf{Erro Silencioso} & \textbf{Acurácia Tática} & \textbf{Latência} \\
\midrule
Regras (Regex) & 0,67 [0,42--0,84] & 27\% [7\%--50\%] & 33\% [9\%--59\%] & 67\% [50\%--83\%] & 0,000 s \\
Laya Multilíngue (CPU) & 0,26 [0,11--0,41] & 47\% [22\%--72\%] & 53\% [27\%--80\%] & 37\% [20\%--53\%] & 1,114 s \\
Ablação: Sem regra <15c & 0,26 [0,11--0,41] & 47\% [22\%--72\%] & 53\% [27\%--80\%] & 37\% [20\%--53\%] & 1,160 s \\
Ablação: Base em inglês & 0,52 [0,30--0,68] & 33\% [11\%--57\%] & 20\% [0\%--46\%] & 57\% [40\%--73\%] & 0,329 s \\
Nuvem (GPT-4o-mini) & 0,95 [0,78--1,00] & 0\% [0\%--0\%] & 0\% [0\%--0\%] & 97\% [90\%--100\%] & 0,450 s \\
\bottomrule
\end{tabular}
\end{table}

A heurística de expressões regulares obteve macro-F1 de 0,67, mas incorreu em 27\% de alarmes falsos, rotulando boletins informativos da Defesa Civil como ``pânico moral'' devido a termos como ``ALERTA''. O Laya cru acusou dispersão severa de confiança decorrente de temperaturas descalibradas na saída. A superioridade do modelo em nuvem (macro-F1 de 0,95) estabelece o teto semântico, reforçando que a inferência local preserva a privacidade do cidadão mas exige calibração pós-hoc e limiares seletivos.

\subsection{Avaliação Comportamental Simulada (Experimento 2)}

A Tabela~\ref{tab:resultados_e2} apresenta os resultados do Experimento~2 sob protocolo de simulação controlada com personas automatizadas, classificado honestamente como \textbf{inconclusivo} quanto ao comportamento de voluntários humanos reais.

\begin{table}[t]
\centering
\caption{Resultados do Experimento 2 (8 sessões / 160 decisões simuladas).}\label{tab:resultados_e2}
\scriptsize
\begin{tabular}{lccc}
\toprule
\textbf{Condição Experimental} & \textbf{Compartilhamento Reativo} & \textbf{Tempo Inspeção} & \textbf{Contestação (7a)} \\
\midrule
Condição A: Feed padrão & 63,7\% [55,0\%--72,5\%] & 3,25 s [3,13--3,36 s] & 0,0\% [0,0\%--0,0\%] \\
Condição B: FiltroSoberano & 35,0\% [28,8\%--41,2\%] & 6,25 s [6,06--6,43 s] & 15,6\% [13,7\%--17,5\%] \\
Condição C: Aviso genérico & 47,5\% [42,5\%--53,8\%] & 4,55 s [4,42--4,69 s] & 0,0\% [0,0\%--0,0\%] \\
Condição D: Sem botão 7a & 40,0\% [32,5\%--47,5\%] & 5,57 s [5,39--5,73 s] & 0,0\% [0,0\%--0,0\%] \\
\bottomrule
\end{tabular}
\end{table}

Na simulação, a exibição do rótulo pedagógico produziu redução de 45,1\% na propensão de compartilhamento frente ao controle (de 63,7\% para 35,0\%), quase duplicando o tempo de leitura (de 3,25~s para 6,25~s). A taxa de contestação situou-se em 15,6\% [13,7\%--17,5\%], abaixo do teto de 20\%.

\subsection{Discussão e Limitações}

No padrão metodológico rigoroso: \textit{observamos} que o classificador não calibrado em CPU apresentou macro-F1 de 0,26 e 47\% de falsos positivos, enquanto a interface pedagógica simulada operou com taxa de contestação de 15,6\%; \textit{comparado a} expressões regulares (que falham por cegueira contextual em 27\% dos casos) e a LLMs proprietários em nuvem (macro-F1 de 0,95), \textit{isso sugere} que a detecção de retórica apelativa não pode ser resolvida por regras léxicas estáticas, e que encoders compactos de 322M requerem indispensavelmente calibragem pós-hoc de probabilidades e aceleração via WebGPU; \textit{não permite concluir}, contudo, que leitores humanos reais sob estímulo político reativo responderão com o mesmo distanciamento observado no ambiente simulado, nem que o modelo opere com confiabilidade sem prévio ajuste fino supervisionado no Brasil.

\section{Conclusão}

Este trabalho apresentou o FiltroSoberano, demonstrando a viabilidade de uma arquitetura de Sistema~1 local no navegador para diagnosticar manipulação retórica e pânico moral sem tráfego de rede e sem censura. Os experimentos falsificaram a suposição de que checkpoints genéricos compactos podem operar diretamente em português sem calibração, quantificando o custo da soberania local e fundamentando as diretrizes para acelerar encoders via WebGPU e calibrar probabilidades com garantias de privacidade.

\bibliographystyle{splncs04}
\bibliography{references}

\end{document}
